%%%%%%%%%%%%%%%%%%%%%%%%%%%%%%%%%%%%%%%%%%%%%%%%%%%%%%%%%%%%%%%%%%%%%%
% Lớp 10 | Xác suất | Xác suất của biến cố
% Dạng 0D0?2-8: Liên quan bàn tròn hoặc hóan vị lặp
% Số câu: 1  (sinh tự động bởi tach_cau.py; từ nay thầy sửa trực tiếp file này)
%%%%%%%%%%%%%%%%%%%%%%%%%%%%%%%%%%%%%%%%%%%%%%%%%%%%%%%%%%%%%%%%%%%%%%

%%%%%%%%%%%%%%%% ĐÚNG/SAI %%%%%%%%%%%%%%%%

% nguồn: MaHoaNH_DungSai_Lop10
\begin{ex}%[0D0H2-8]%[ma:5f970e55]
	Một nhóm bạn có $4$  bạn nam và $4$ bạn nữ ngồi ngẫu nhiên vào một bàn tròn. Xác suất để các bạn nam và nữ ngồi xen kẽ nhau bằng
	 \choiceTF
		{$\dfrac{1}{2}$} 
		{$\dfrac{1}{4}$} 
		{\True $\dfrac{1}{35}$} 
		 {$\dfrac{1}{70}$}
	 \loigiai{
		Số phần tử không gian mẫu  $n\left( \Omega \right)=7!$.\\
		Gọi $A$ là biến cố:  “Các bạn nam và nữ ngồi xen kẽ”.
		Khi đó ta có $n(A)=3! \cdot 4!$.  \\
	Vậy $\mathrm{P}(A)=\dfrac{n(A)}{n\left(  \Omega \right)}=\dfrac{3! \cdot 4!}{7!}=\dfrac{1}{35}$.}
\end{ex}
